\subsection{FRM Exam 2005: Monthly VaR interpretation (95\%)}

\textbf{Enunciado}

\emph{The 10-Q report of ABC Bank states that the monthly
VaR of ABC Bank is USD 10 million at the 95\% confidence level. What is
the proper interpretation of this statement?}
\begin{enumerate}[label=\alph*.]
\item {If we collect 100 monthly gain/loss data of ABC Bank, we will always see five months with losses larger than \$10 million.}
\item \hl{There is a 95\% probability that the bank will lose less than \$10 million over a month.}
\item {There is a 5\% probability that the bank will gain less than \$10 million each month.}
\item {There is a 5\% probability that the bank will lose less than \$10 million over a month.}
\end{enumerate}

\subsubsection*{Solución}


\(\VaR_{0.95}(L)=10\)M significa:
\[
\Prob(L\le 10\text{M})\approx 0.95,\qquad \Prob(L>10\text{M})\approx 0.05.
\]
Eso es exactamente el inciso (b). El inciso (a) es falso por el ``always'' (el conteo es aleatorio), (c) mezcla el signo y (d) invierte 95\% y 5\%.











